% ============================================================
%  Chapter 2 — Websites and Browsers
%  The Secret Life of the Internet
% ============================================================

\chapter{Websites and Browsers}

\chapteropening{%
  There are billions of websites on the internet.
  How do you get to exactly the right one?
  Let Bee Browser be your guide!
}
\chapterbadge{Buzzing Through Website City}

% --- Opening story ---
\section*{Welcome to the Website Neighbourhood}

Dot had guided Sam to the edge of the internet highway.
Now it was time to visit a website.

A cheerful yellow bee wearing tiny goggles buzzed into view.

\character{Bee Browser}{Hello! I'm Bee Browser.
  My job is to take you anywhere on the internet you want to go.
  Just tell me the address and off we fly!}

Sam looked around.
In every direction there were colourful buildings —
a towering library, a bright game arcade, a cosy weather station,
a giant video cinema, and hundreds more stretching to the horizon.

\begin{picturethis}
  Imagine a huge colourful city where every single building is a website.
  The library building is a news website.
  The arcade is a games website.
  The cinema is a video website.
  Each building has its very own address — and Bee Browser knows how to find them all!
\end{picturethis}

\sectionrule

% --- What is a website? ---
\section*{What Is a Website?}

\begin{funfact}
There are over 1.1 billion websites on the internet, but fewer than 200 million are actively visited.
\end{funfact}

A \term{website} is a collection of pages stored on a special computer called a \term{server}.
Those pages can hold words, pictures, videos, buttons, sounds, and much more.

Every website has a unique address called a \term{URL}
(which stands for Uniform Resource Locator — quite a mouthful, so most people just say ``web address'').

A web address looks like this:

\begin{center}
  \large\ttfamily\color{InternetBlue}
  https://www.example.com
\end{center}

\begin{quickmap}
  \textbf{Breaking down a web address:}
  \begin{enumerate}[label=\textbf{\arabic*.}]
    \item \texttt{https://} — a signal that the connection is \emph{secure} (the ``s'' stands for secure).
    \item \texttt{www} — short for ``World Wide Web''.
    \item \texttt{example} — the name of the website.
    \item \texttt{.com} — the type of website (commercial, organisation, country, etc.).
  \end{enumerate}
\end{quickmap}

\begin{diagram}[What happens in a flash when you press Enter.]
  \begin{tikzpicture}[>=Stealth, node distance=0.8cm]
    \node[draw=InternetBlue, rounded corners, fill=SkyBlue!25, minimum width=3.6cm, minimum height=0.9cm] (s1) {You type URL};
    \node[draw=WarmOrange!80!black, rounded corners, fill=WarmOrange!22, minimum width=3.6cm, minimum height=0.9cm, below=of s1] (s2) {Browser asks DNS};
    \node[draw=LeafGreen!80!black, rounded corners, fill=LeafGreen!20, minimum width=3.6cm, minimum height=0.9cm, below=of s2] (s3) {DNS returns IP};
    \node[draw=SoftPurple!80!black, rounded corners, fill=SoftPurple!20, minimum width=3.6cm, minimum height=0.9cm, below=of s3] (s4) {Browser contacts Server};
    \node[draw=CoralRed!80!black, rounded corners, fill=CoralRed!16, minimum width=3.6cm, minimum height=0.9cm, below=of s4] (s5) {Server sends page};
    \node[draw=InternetBlue!80!black, rounded corners, fill=SunYellow!25, minimum width=3.6cm, minimum height=0.9cm, below=of s5] (s6) {Browser displays page};
    \draw[thick,->,InternetBlue] (s1) -- (s2);
    \draw[thick,->,InternetBlue] (s2) -- (s3);
    \draw[thick,->,InternetBlue] (s3) -- (s4);
    \draw[thick,->,InternetBlue] (s4) -- (s5);
    \draw[thick,->,InternetBlue] (s5) -- (s6);
  \end{tikzpicture}
\end{diagram}

\sectionrule

% --- What is a browser? ---
\section*{What Is a Browser?}

A \term{browser} is a programme on your device that lets you visit and view websites.

Think of it this way:

\begin{picturethis}
  Bee Browser is like a taxi driver.
  You tell the taxi driver the address — say, \texttt{www.mycoollibrary.com} —
  and the taxi driver takes you straight there.
  The browser is the taxi.
  The web address is where you want to go.
  The website that appears on your screen is your destination.
\end{picturethis}

Popular browsers you might recognise:
\begin{itemize}
  \item Chrome
  \item Firefox
  \item Safari
  \item Edge
\end{itemize}

They all do the same basic job: \textbf{take your web address and bring the website back to your screen}.

\sectionrule

% --- How a browser fetches a page ---
\section*{What Happens When You Press Enter?}

It feels like magic, but here is what really happens in just a fraction of a second:

\begin{quickmap}
  \begin{enumerate}[label=\textbf{\arabic*.}]
    \item You type a web address and press Enter.
    \item Your browser asks a special helper called a \term{DNS server},
          ``What is the number address for this website?''
    \item The DNS server replies with the server's \term{IP address} (a number like 192.0.2.1).
    \item Your browser travels to that IP address and says, ``Please send me your page!''
    \item The server sends back the page in pieces.
    \item Your browser puts all the pieces together and displays them on your screen.
  \end{enumerate}
\end{quickmap}

\character{Bee Browser}{It's like using a phone book!
  You look up a name, find the phone number, and call.
  DNS is the internet's phone book.}

\sectionrule

% --- HTML teaser ---
\section*{What Are Websites Made Of?}

Websites are built with a special language called \term{HTML}
(HyperText Markup Language).
HTML is a set of instructions that tells your browser what to show:
headings, paragraphs, pictures, links, and buttons.

You do not need to know HTML to use the internet.
But it is fun to know that every website you visit is built from
these hidden instructions — like a recipe that your browser reads and cooks into the page you see.

\sectionrule

\begin{newwords}
  \begin{description}[labelwidth=3.2cm, leftmargin=3.4cm]
    \item[\term{Website}]  A collection of pages on the internet, stored on a server.
    \item[\term{Browser}]  A programme (like Chrome or Safari) that lets you visit websites.
    \item[\term{URL}]      A web address — the unique location of a website.
    \item[\term{DNS}]      A system that turns website names into number addresses.
    \item[\term{IP address}] A number that identifies a device or server on the internet.
    \item[\term{HTML}]     The language used to build web pages.
  \end{description}
\end{newwords}

\sectionrule

\begin{tryit}
  \textbf{Match the address to the building!}

  Draw lines to match each website type to the right building:

  \begin{center}
    \begin{tabular}{ll}
      \textbf{Website type} & \textbf{Building in the city} \\[4pt]
      News website          & Library \\
      Video website         & Cinema \\
      Games website         & Arcade \\
      Weather website       & Weather station \\
      Shopping website      & Shopping mall \\
    \end{tabular}
  \end{center}

  \emph{Now try opening a browser on your device and type in a web address you know.
  What building would that website be in our internet city?}
\end{tryit}

\sectionrule

\begin{bigidea}
  A \textbf{website} is a place on the internet with its own address.
  A \textbf{browser} is like a taxi that takes you there —
  just type the address and it fetches the page for you.
\end{bigidea}

\character{Dot}{Now you know how to visit websites!
  But wait — how does your tablet even connect to the internet in the first place?
  Let's visit Willa Wi-Fi and Rory Router next!}
